\chapter{지휘자 없는 기계}
% full version: chapters/02_glutcomputer.tex

여러분의 컴퓨터에는 지휘자가 있다. 바로 클록 발생기이다. 클록은 1초에 수십억 번 박자를 주고, 회로의 모든 소자는 그 박자에 맞춰 한꺼번에 상태를 바꾼다. 뇌에는 이런 지휘자가 없다. 뉴런은 저마다 따로 일한다. 신호는 아무 때나 들어오고 아무 때나 나간다. 이런 기계에서 \nt{GLUT} 뉴런, 즉 “플러스”만으로 무엇을 계산할 수 있는지 살펴보자.

\section*{구멍 난 양동이}

뉴런은 구멍 난 양동이로 생각하면 편하다. 들어오는 스파이크 하나하나가 물 한 바가지이다. 물은 구멍으로 조금씩 새어 나간다. 바가지가 자주 들어오면 물이 가장자리까지 차오르고, 양동이가 뒤집힌다. 뉴런이 딸깍하고, 양동이는 비고, 모든 것이 처음부터 다시 시작된다. 바가지가 드물게 들어오면 그사이에 물이 새어 나가 버려 아무 일도 일어나지 않는다.

\pencilfig{2}{\pencilart{2}}{뉴런은 구멍 난 양동이이다. 입력 방울이 양동이를 채우고 물은 새어 나간다. 물이 선에 닿으면 딸깍.}

구멍이야말로 뉴런이 논리 게이트와 다른 핵심이다. 뉴런은 들어온 스파이크를 영원히 기억하지 않고 몇 밀리초 동안만 기억한다. 이 창 안에 들어온 것은 모두 더해지고, 그보다 앞선 것은 잊힌다.

\section*{“또는”과 “그리고”}

바가지 하나로 양동이가 뒤집힌다면 뉴런은 어떤 입력에든 반응한다. 이것이 “또는”이다. 바가지 두 개가 필요하면 흥미로워진다. 지휘자가 없으니 “둘 다 왔는가?”라는 물음은 \emph{얼마의 시간 안에}라는 조건이 붙기 전에는 의미가 없다. 뉴런은 첫 바가지가 새어 나가기 전에, 즉 두 스파이크가 거의 동시에 도착해야만 발화한다. 이것이 “시간 속의 그리고”, 곧 동시성 검출기이다. 어떤 뉴런의 창은 몇 밀리초이고, 청각 뉴런의 창은 1밀리초보다도 짧다.

\section*{“플러스”만으로는 못 하는 일}

“아니다”를 만들어 보자. 입력이 조용할 때 일하는 뉴런이다. 되지 않는다. 침묵은 양동이에 물을 붓지 않는다. 뉴런 하나만이 아니라 “플러스”로만 된 어떤 네트워크도 마찬가지이다. 입력을 세게 하면 그 네트워크의 어떤 뉴런도 덜 일하게 되지 않는다. 여기서 가장 곤란한 성질이 나온다. \emph{활동을 활동으로 끌 수는 없다}. 번진 불을 멈추는 방법은 하나뿐이다. 장작을 더 넣지 않는 것이다.

자연은 우회로를 찾았다. 여러 감각 기관에서는 같은 자극을 두 무리의 뉴런이 전한다. 눈에서는 어떤 세포가 “밝아졌다”를, 다른 세포가 “어두워졌다”를 알린다. 전선 하나 대신 “예”와 “아니요” 두 가닥의 전선 쌍을 쓰면 “아니다”는 공짜로 얻는다. 두 전선을 바꿔 끼우기만 하면 된다. 이런 쌍을 쓰면 “플러스”만으로 어떤 논리 회로든 만들 수 있고, 계산이 끝났다는 것까지 알 수 있다. 모든 쌍에서 정확히 한 가닥에 불이 들어오면 끝난 것이다. 엔지니어들은 이런 방식으로 어떤 지연에서도 올바르게 동작하는 비동기 회로를 만든다.

\section*{지연으로 만드는 시간}

지휘자가 없으면 공통 시계도 없지만, 시간은 이렇게도 계산할 수 있다. 옆에서 나는 소리는 가까운 귀에 먼 귀보다 조금 먼저, 1밀리초보다 짧은 차이로 도착한다. 왼쪽 귀에서 나온 전선을 동시성 검출기 줄을 따라 왼쪽에서 오른쪽으로 깔고, 오른쪽 귀에서 나온 전선은 오른쪽에서 왼쪽으로 깐다. 두 신호는 서로를 향해 달려가 둘이 동시에 도착하는 지점에서 만난다. 바로 그 자리의 검출기가 발화한다. 시간 차가 발화한 뉴런의 \emph{번호}로, 즉 소리가 오는 방향으로 바뀐 것이다. 정밀도는 10마이크로초 정도이다. 뉴런 하나하나는 훨씬 덜 정밀하지만, 검출기가 많다는 것이 이를 메워 준다.

\pencilfig{102}{%
% delay line: the left-ear wire runs right along the top, the right-ear wire runs left along the bottom
\draw[pen=pcBlue] (0,0.6) -- (5.6,0.6); \draw[pen=pcBlue] (6.0,-0.6) -- (0.4,-0.6);
\foreach \i in {1,...,7}{
  \pgfmathsetmacro{\x}{0.75*\i}
  \draw[pen=pcBlue] (\x,0.6) -- (\x,0.24); \draw[pen=pcBlue] (\x,-0.6) -- (\x,-0.24);
  \ifnum\i=5 \draw[dotpen=pcAmber, wash=pcAmber, line width=1.1pt] (\x,0) circle (0.2);
  \else \draw[dotpen=pcBlue, wash=pcBlue] (\x,0) circle (0.2); \fi}
\draw[arr=pcBlue] (0.95,0.78) -- (1.75,0.78); \draw[arr=pcBlue] (5.3,-0.78) -- (4.5,-0.78);
\node[lbl, anchor=east] at (-0.05,0.6) {왼쪽 귀};
\node[lbl, anchor=west] at (6.05,-0.6) {오른쪽 귀};
\node[lbl, text=pcAmber!70!black, anchor=south] at (3.75,0.95) {여기서 만났다};
\draw[arr=pcAmber!70!black] (3.75,0.95) -- (3.75,0.68);
% sound source on the left
\draw[penlite] (-1.25,-0.25) -- (-1.05,-0.25) -- (-0.8,-0.45) -- (-0.8,0.05) -- (-1.05,-0.15) -- (-1.25,-0.15) -- cycle;
\foreach \r in {0.2,0.35}{\draw[penlite] ($(-0.75,-0.2)+(-40:\r)$) arc (-40:40:\r);}
\node[lbl, anchor=north] at (-0.95,-0.5) {왼쪽 소리};
}{왼쪽에서 난 소리는 왼쪽 귀에 먼저 닿는다. 그 신호가 더 멀리 달려가므로 만나는 곳은 가운데보다 오른쪽이다. 발화한 뉴런의 번호가 곧 소리의 방향이다.}

\section*{기억은 모닥불이다}

서로를 강하게 흥분시키는 뉴런 무리를 생각하자. 짧은 신호로 불을 붙이면 이 무리는 스스로 장작을 넣는 모닥불처럼 제 활동을 유지한다. 신호는 끝났는데 무리는 계속 타오른다. 신호가 있었다는 것을 기억하는 셈이다. 이것이 메모리 셀이다. 신호가 너무 짧으면 모닥불은 미처 타오르지 못하고 꺼진다.

그런데 문제가 있다. 이런 기억은 어떻게 지우는가? 활동으로는 방법이 없다. 뉴런이 지쳐 모닥불이 다 타 버릴 때까지 기다리는 수밖에 없다.

\section*{사건이 시작하는 타이머}

무리들을 사슬로 잇자. 첫 무리가 둘째에 불을 붙이고, 둘째가 셋째에 붙인다. 첫 무리를 밀면 도미노가 쓰러지듯 사슬을 따라 파동이 달린다. 쓰러진 도미노의 번호가 밀고 나서 시간이 얼마나 지났는지 알려 준다. 이것은 늘 째깍거리는 시계가 아니라, 사건이 켜고 한 번만 도는 타이머이다. 노래하는 새가 노래할 때 한 영역의 뉴런들은 엄격히 차례대로 발화한다. 이런 사슬과 아주 비슷하다.

이렇게 “플러스”만으로도 많은 것을 만들 수 있다. 모자란 것은 세 가지이다. 여럿 중 하나를 고르는 일, 명령에 따라 기억을 지우는 일, 네트워크가 불길에 휩싸이지 않게 붙잡는 일이다. 이 셋은 모두 “마이너스”가 해 준다. 다음 장의 주제이다.

\fullversion{2}
